% =====================================================================
%  Chapter 7  ·  Session 7 --- Hardening and access control
% =====================================================================
\chapter{Hardening and access control}

The firewall of Session~4 decides where traffic may come from, but it says
nothing about who is on the other end or what they may do once connected.
This chapter covers authentication (proving who one is), authorisation
(deciding what one may do) and least privilege. These questions weigh more
in the cloud than elsewhere, because the whole infrastructure is managed
through a web console and an API reachable from the Internet, so the account
that controls it is the way in.

\section{Identity as the new perimeter}

The CSA Security Guidance calls identity and access management the new
perimeter of the cloud, and the reason is this consolidation: a provider
gathers all administration into one web console and one API, and both are
reachable from the Internet (CSA, 2024, pp.~103--104). Often nothing more
than a username and a password stands in front of them, and most breaches of
cloud-native systems accordingly begin with a failure of identity management
rather than with a network attack (CSA, 2024, p.~104). A stolen
administrator credential is therefore no longer a foothold inside a building
that must still be crossed; it can be control of everything at once. The
perimeter formed by IAM covers authentication as well as authorisation,
every kind of actor (people, systems and code) and every point an attacker
can reach, from a phished password to a key left in a published container
build file (CSA, 2024, p.~116). Five terms supply the vocabulary (CSA, 2024,
pp.~105--106).

\begin{definition}[Core IAM terms]\label{def:iam}
An \textbf{entity} describes a unique actor: a person, device, application or
system. An \textbf{identity} describes the entity's verifiable expression
within a namespace; one entity may hold several. \textbf{Authentication}
describes the process of verifying an identity. \textbf{Authorisation}
describes the decision to permit or deny an authenticated identity access to
a resource. An \textbf{entitlement} describes the mapping of identities to
authorisations, often with conditions (identity~X may access resource~Y when
attributes~Z hold); a table of these is an \emph{entitlement matrix}.
\end{definition}

Authentication and authorisation must be kept apart, because a system can
get the first right and the second wrong: it identifies a user correctly and
then grants far more access than the task needs. Least privilege exists to
prevent that second failure, and the entitlement matrix, written down and
reviewed, is the document that shows whether it holds.

\section{Authentication and multi-factor authentication}

Authentication rests on \emph{factors}, that is, categories of evidence for
a claim of identity. There are three, and combining them is the most
effective identity control available.

\begin{definition}[Multi-factor authentication]\label{def:mfa}
\textbf{Multi-factor authentication} (MFA) describes the requirement of
evidence from more than one independent category: \emph{something one knows}
(a password or PIN), \emph{something one has} (a phone, hardware token or
key) and \emph{something one is} (a biometric). A stolen password alone is
then insufficient, because it satisfies only one factor.
\end{definition}

The cloud raises the need for MFA for two reasons (CSA, 2024, pp.~112--113).
The first is \emph{broad network access}, the second essential
characteristic of Session~1: cloud services are reached over the Internet,
so a lost or stolen credential can be used from anywhere, and the attacker
no longer needs to be on the local network. The second is single sign-on
across federated services, because one compromised set of credentials then
opens many services at once. The Guidance accordingly ranks MFA among the
most effective measures against account takeover, and a cloud account
guarded by a password alone is a high risk (CSA, 2024, p.~113). The factors
differ in strength: hardware tokens and phishing-resistant methods such as
FIDO passkeys are stronger than a code read from an app, which in turn is
stronger than a password alone (CSA, 2024, pp.~113--114). MFA is also the
basis of \emph{conditional access}, in which the system weighs the
circumstances of a login, such as an unknown device, an unusual country, or
the impossible travel of two logins from distant places minutes apart, and
demands more assurance, or refuses, when they look wrong (CSA, 2024,
pp.~106, 116).

\begin{realworld}{Colonial Pipeline, 2021}
The largest fuel pipeline in the United States was shut down for days,
causing shortages along the East Coast, after a ransomware crew entered
through a single account. The account belonged to a legacy virtual private
network, its password had appeared in a batch of leaked credentials
elsewhere, and it had no multi-factor authentication. One reused password
without a second factor was sufficient. With MFA the stolen password alone
would have failed (Definition~\ref{def:mfa}). MFA on every account that
reaches inside is a baseline, not advice.
\end{realworld}

\section{Authorisation: RBAC, PBAC and ABAC}

Once an identity is authenticated, authorisation decides what it may do, and
three models express that decision, in rising order of flexibility (CSA,
2024, pp.~105, 118--119).

\begin{definition}[Authorisation models]\label{def:authz}
\textbf{Role-based access control} (RBAC) describes granting permissions to
\emph{roles} and assigning identities to roles (administrator,
member). \textbf{Policy-based access control} (PBAC) describes expressing
permissions as explicit policies evaluated at access time.
\textbf{Attribute-based access control} (ABAC) describes deciding from
\emph{attributes} of the identity, resource and context (department, resource
sensitivity, time, location, MFA status), which enables fine-grained,
context-aware decisions.
\end{definition}

RBAC is simple and familiar, and for a small service it is often sufficient.
Its weakness is coarseness, because everyone in a role receives the same
access regardless of circumstance. ABAC and PBAC add that circumstance:
access only from a managed device, only during working hours, only if MFA
was used. The Guidance therefore names PBAC supporting ABAC as the preferred
model and recommends ABAC and PBAC over RBAC where the platform supports
them (CSA, 2024, pp.~115, 122). Table~\ref{tab:authz} compares the three.

\begin{table}[htbp]
\centering
\small
\begin{tabular}{@{}l p{3.5cm} p{3.0cm} p{3.3cm}@{}}
\toprule
\textbf{Model} & \textbf{Decides from} & \textbf{Strength} & \textbf{Weakness} \\
\midrule
RBAC & the identity's \emph{role} & simple, familiar, easy to audit & coarse: all in a role get the same access \\
PBAC & explicit \emph{policies} evaluated at access time & central and expressive & policies must be written and maintained \\
ABAC & \emph{attributes} of identity, resource and context & fine-grained, context-aware & complex to design and reason about \\
\bottomrule
\end{tabular}
\caption{Three authorisation models. All three express who may do what;
they differ in what the decision is based on and in the granularity bought at
the cost of complexity. The Guidance prefers ABAC/PBAC where the platform
supports them.}
\label{tab:authz}
\end{table}

Whichever model is used, one principle governs it.

\begin{definition}[Least privilege]\label{def:leastpriv}
The principle of \textbf{least privilege} describes the rule that every
identity is granted the minimum access its task requires, and no more. Excess
privilege is a hidden liability rather than a convenience: it is the
additional damage available to any mistake and to any attacker who takes the
identity over.
\end{definition}

Least privilege has already appeared twice under other names: as
\texttt{root} versus \texttt{sudo} in Session~2, where administrative power
is not held while it is not in use, and as the firewall's default-deny in
Session~4, where no network access is granted that has not been justified.
Here it becomes the shape of every authorisation decision.

\section{Protection domains and the access control matrix}

RBAC, ABAC and least privilege are the modern form of a model from
operating-systems theory (Tanenbaum \& Bos, 2015, pp.~602--604). A system
holds \emph{objects} to be protected (files, devices, processes, memory),
each with a name and a fixed set of operations. A \emph{protection domain}
is a set of \emph{(object, rights)} pairs: for one actor, which operations
are permitted on which objects. With every domain as a row and every object
as a column, the result is the \emph{access control matrix}, the most
general statement of who may do what to what. Every authorisation scheme in
this chapter compresses that matrix so that it can be managed at scale. RBAC
groups rows that share a pattern into roles. ABAC computes a cell's contents
from attributes at access time rather than storing it. An entitlement matrix
(Definition~\ref{def:iam}) is the access control matrix written down for
review. Least privilege, lastly, is Tanenbaum and Bos's own rule for filling
the matrix, the \emph{Principle of Least Authority} (POLA), also called need
to know: each domain holds the minimum objects and rights needed for its
work.

The Linux machine of Session~2 already works this way. In UNIX and Linux a
process's protection domain is defined by its \emph{user and group
identifiers} (UID and GID): at login the shell receives them from the
password database, every process started from it inherits them, and the
pair determines which files and devices may be read, written or executed.
Changing the pair changes the domain, and two mechanisms for doing so carry
into the cloud. First, each process has a user half and a kernel half, and a
\emph{system call} switches between them. Because the kernel half can touch
physical memory and every disk that the user half cannot, a system call is a
domain switch to a more privileged domain, the same trap that made
virtualisation work in Session~2. Secondly, running a program marked
\emph{set-UID} gives it a different effective identity for the duration.
This is the mechanism behind \texttt{sudo}, and the reason \texttt{sudo} is
a privilege escalation to be granted sparingly.

\section{Single sign-on and privileged access}

Zero Trust validates every request, and in a service built from many
microservices a naive implementation would therefore prompt the user for a
password at every hop. Comer (2021) gives the example of one task, finding a
colleague's email address, that involves two services and asks the user to
log in at each step (Comer, 2021, pp.~237--238). The situation is worse when
every service keeps its own user list and passwords, because each service
then grants privileges on its own, with no coordination. The answer is a
central \emph{identity management} (IdM) system offering \emph{single
sign-on} (SSO) (Comer, 2021, p.~238). The user has one set of credentials,
held in one place together with their access rights, every service consults
that system, and the system can return a digital capability that other
services accept, so the user enters credentials once per task. SSO is
therefore what makes per-request checking workable, and it is the cloud
form of the \emph{federation} at the centre of the Guidance's identity
domain (CSA, 2024, pp.~107--109). In Azure the central system is Microsoft
Entra ID, which is why accounts, MFA and conditional access are managed
there rather than in each application.

The same centralisation enables a control for the accounts that matter
most: \emph{privileged access management} (PAM), identity management for the
administrators themselves (Comer, 2021, pp.~238--239). There is no single
master password that opens every system; each member of staff holds
administrative rights only on the systems that person is responsible for.
The PAM system logs all accesses and failed login attempts, and some systems
alert a manager when a privileged identity is used in a suspicious way, for
example to reach a system for which it holds only secondary responsibility,
or several such systems within a short time. The Guidance treats the same
topic under privileged user management and recommends the strongest
available authentication and monitoring for these accounts (CSA, 2024,
pp.~117--118). PAM is thus least privilege applied to the administrators
themselves.

\section{Credentials and the management plane}

Two practical hazards remain. The first is \emph{static credentials}:
long-lived secrets such as API keys embedded in code or configuration,
which are a permanent gift to an attacker who reads the repository. The
Guidance recommends eliminating them as far as possible (CSA, 2024,
pp.~119--120), and where a human needs elevated access, granting it
\emph{just in time} (JIT): requested for a session, approved, and revoked
when the session ends, so that no permanent key exists to steal (CSA, 2024,
p.~105). The second hazard is the management plane itself. Because it
controls everything, administrative access to it gets the strongest
treatment available (CSA, 2024, pp.~116--118): MFA without exception, least
privilege enforced, and administrative logins restricted by location where
possible. Session~9 adds that every action on it is logged and watched.

\begin{example}[Hardening the Nextcloud administrator]\label{ex:mfa-trc}
Nextcloud has one all-powerful administrator account. \emph{Threat}: the administrator's password is phished or reused
from a breach elsewhere, and an attacker logs in with full control of every
user's files. \emph{Risk}: the likelihood is real, because password reuse and
phishing are everywhere, and the impact is maximal, because the account can
read, alter and delete everything, affecting all three CIA properties at once.
\emph{Controls}, layered: MFA on the administrator account, so that a
stolen password alone fails (Definition~\ref{def:mfa}); only the access they
need for ordinary members, not administrative rights (least privilege); and
no single shared administrator login, so that each action can be traced to
a person, which is also what makes the logs of Session~10 meaningful. No
single control removes the threat, but together they raise its cost.
\end{example}

\paragraph{Reading for this chapter.} CSA Security Guidance, Domain~5, in
particular \S5.1 (How IAM is Different in the Cloud, and Fundamental Terms,
pp.~104--106), \S5.2 (Federation, pp.~107--111), \S5.3 (Strong Authentication
\& Authorization, pp.~112--118), \S5.4 (IAM Policy Types, pp.~118--119) and
\S5.5 (Least Privilege \& Automation, pp.~119--121); the domain runs
pp.~103--123. Comer, Chapter~17, \S\S17.5--17.7 (Zero Trust, Identity
Management and single sign-on, and Privileged Access Management,
pp.~237--239); Tanenbaum and Bos, \emph{Modern Operating Systems}, \S9.3
(Protection Domains and the access control matrix, and their UNIX/Linux
realisation in UID/GID, pp.~602--604). Optional deepening: J\o sang,
Chapter~8 (User Authentication, pp.~165--190) and Chapter~9 (IAM---Identity
and Access Management, pp.~191--214). The Microsoft Learn module is \emph{Describe Azure identity}
(Microsoft Entra ID, MFA and conditional access); students enable MFA on their
own account.

\section{Summary}

When a cloud is managed through Internet-facing consoles and APIs, identity
is the perimeter (Definition~\ref{def:iam}), and its two questions are kept
apart: authentication proves who the requester is, authorisation decides
what the requester may do. Because broad network access turns a stolen
credential into an easy account takeover, multi-factor authentication,
something one knows, has and is (Definition~\ref{def:mfa}), together with
conditional access, is the baseline. Authorisation is expressed through
RBAC, PBAC or ABAC (Definition~\ref{def:authz}), the context-aware models
being preferred where available, and all of them serve least privilege
(Definition~\ref{def:leastpriv}), which is grounded in the access control
matrix of operating-systems theory. In practice this means replacing static
credentials by just-in-time access and giving the management plane the
strongest controls. On the red thread, the Nextcloud administrator gains MFA
and members lose privileges they never needed (Example~\ref{ex:mfa-trc}), so
every action can be traced to a person, which the monitoring and logs of
Sessions~9 and~10 depend on. Session~8 turns to the application itself.

\section*{Self-check}
\addcontentsline{toc}{section}{Self-check}

Sketch answers follow the exercises.

\begin{exercise}[Authentication vs authorisation]
Distinguish authentication from authorisation, and give an example of a system
that gets the first right and the second dangerously wrong.
\end{exercise}

\begin{exercise}[The three authentication factors]
Name the three authentication factors, give an example of each, and explain
why MFA defeats a stolen password.
\end{exercise}

\begin{exercise}[Why the cloud needs MFA]
Why does the CSA Guidance hold that the cloud increases the need for strong
MFA? Tie the answer to one of the NIST essential characteristics from
Session~1.
\end{exercise}

\begin{exercise}[RBAC vs ABAC]
Contrast RBAC and ABAC, and state which the Guidance prefers and why. How does
each relate to least privilege?
\end{exercise}

\begin{exercise}[Shared admin account as threat--risk--control]
Write, in the vocabulary of Definition~\ref{def:trc}, the threat, risk and
controls for a single shared Nextcloud administrator account protected only by
a password.
\end{exercise}

\paragraph{Sketch answers.}
{\small
(1)~Authentication verifies identity; authorisation decides permitted
actions. A system authenticates a user correctly but then grants them
administrator rights they do not need: correct identity, excessive
authorisation. (2)~Something one knows (password/PIN), something one has
(phone/hardware token), something one is (biometric); MFA requires more than
one category, so a stolen password satisfies only one factor and fails alone.
(3)~Broad network access: cloud services are reached over the Internet, so a
lost or stolen credential leads to account takeover without the attacker
being anywhere near the victim, and passwords alone cannot carry that risk.
(4)~RBAC grants permissions to roles and assigns identities to roles, simple
but coarse; ABAC decides from attributes and context, fine-grained and
context-aware. The Guidance prefers ABAC (and PBAC) where supported; both
serve least privilege by granting only what the task and context require.
(5)~Threat: the shared password is phished or reused, giving an attacker full
control of all files. Risk: high likelihood, maximal impact (all three CIA
properties). Controls: require MFA on the admin account, apply least privilege
to members, and replace the shared login with per-person accounts so that
each action can be traced to a person.}


